

\section{Contextualização da Instituição}
~ % Utilizado para saltar linha com recuo
 
Nas primeiras décadas após a Proclamação da República, o processo de industrialização brasileiro intensificou a necessidade de formação de mão de obra qualificada para operar novas tecnologias e métodos produtivos. A demanda por capacitação técnica tornou-se evidente, conduzindo a nação à criação de mecanismos institucionais voltados ao desenvolvimento de competências profissionais.

Nesse contexto, em setembro de 1909, o Decreto nº 7.566, sob a presidência de Nilo Peçanha, instituiu as Escolas de Aprendizes Artífices, sediadas nas capitais dos estados, com oferta de ensino primário profissionalizante e de caráter gratuito. No Ceará, a respectiva unidade foi instalada em 24 de maio de 1910, na Avenida Alberto Nepomuceno, local onde hoje se situa a Secretaria Estadual da Fazenda.

A reorganização administrativa de 1930 vinculou essas escolas ao recém-criado Ministério da Educação e Saúde Pública. Posteriormente, a Lei nº 378, de 1937, conduzida pelo Ministro Gustavo Capanema, transformou as Escolas de Aprendizes Artífices em Liceus Profissionais, resultando, no Ceará, na criação do Liceu Industrial de Fortaleza.

A partir da Segunda Guerra Mundial, a redução das importações de bens industrializados levou o país a estimular o desenvolvimento de sua própria base produtiva, ampliando a demanda por formação técnica especializada. Consequentemente, o Decreto nº 4.121, de 1942, denominou a instituição como Escola Industrial de Fortaleza. No mesmo ano, a Lei Orgânica do Ensino Industrial estruturou o ensino profissional em nível equivalente ao secundário, consolidando a formação técnica como etapa formal de educação de segundo grau.

Em 1953, a Lei nº 3.552 conferiu à instituição o nome de Escola Técnica Federal do Ceará, consolidando-a como referência estadual na formação técnica profissional. Nas décadas seguintes, sua atuação foi ampliada com a criação das Unidades Descentralizadas de Ensino (UNEDs) em Cedro e Juazeiro do Norte.

Com a evolução das demandas tecnológicas do país, a Lei nº 8.948, de 1994, instituiu os Centros Federais de Educação Tecnológica (CEFETs), atribuindo-lhes novas competências, como a oferta de cursos superiores, pós-graduação, pesquisa aplicada e formação de docentes e especialistas. Nesse momento, a instituição passou a denominar-se Centro Federal de Educação Tecnológica do Ceará (CEFET-CE).

A transformação institucional mais abrangente ocorreu com a promulgação da Lei nº 11.892, de 2008, que criou os Institutos Federais de Educação, Ciência e Tecnologia, incorporando os CEFETs a uma nova estrutura multicampi, dotada de autonomia administrativa e pedagógica, com atuação integrada em ensino, pesquisa e extensão. Assim, constituiu-se o Instituto Federal de Educação, Ciência e Tecnologia do Ceará (IFCE), integrante da Rede Federal com áreas de ensino distribuídas pelo estado do Ceará, e comprometidos com uma formação pública, ética, socialmente referenciada e voltada ao desenvolvimento regional. Atualmente o IFCE é composto por 33 \textit{campi} e um polo de inovação; contando ainda com a projeção de expansão de outros 6 \textit{campi}, conforme \citeonline{MEC_IFCE_2024}. 

\subsection{Breve Histórico do Campus Maracanaú}
~

A Lei nº 11.892/2008, ao instituir os Institutos Federais, consolidou a rede de \textit{campi} em todo o território nacional. À época, o IFCE já incluía o \textit{Campus} Maracanaú, inicialmente criado como CEFET–UNED Maracanaú, em 2006.

Localizado na Região Metropolitana de Fortaleza, o município configura-se como o maior polo industrial do Ceará, reunindo setores metalmecânico, alimentício, têxtil, químico e de tecnologia, com aproximadamente {\vm 18.000} empresas instaladas, segundo estimativas governamentais, o que o posiciona como um dos principais eixos produtivos do Estado.

O IFCE – \textit{Campus} Maracanaú foi concebido para atender à crescente demanda por profissionais qualificados, alinhando sua oferta formativa às características do arranjo produtivo local. Seu primeiro curso ofertado, ainda em 2006, foi o curso técnico em Desenvolvimento de \textit{Software}. Atualmente, o \textit{campus} oferta cursos técnicos e superiores, além do Mestrado em Energias Renováveis e da participação no Mestrado Acadêmico em Ciência da Computação, sediado no \textit{Campus} Fortaleza.

Nesse cenário, o Eixo Tecnológico da Indústria, com o apoio da Direção-Geral do \textit{campus}, instituiu, por meio das Portarias \PORTARIACOMISSAO (ANEXO \ref{anexoPortariaComissao}) e \PORTARIACOMISSAOALT (ANEXO \ref{anexoPortariaComissaoAlt}), a comissão responsável pela implantação do \NOMECURSO, reconhecendo-o como formação estratégica para os setores industrial e de serviços.

